\section{Problem Formulation}

Let a GUI visual-search example consist of a screenshot $I$, a target cue $q$, and a human or model scanpath $S = (f_1,\ldots,f_T)$. Standard target-present evaluation assumes that $q$ corresponds to an element in $I$ and measures whether $S$ or the final prediction approaches the ground-truth target.

We extend the output space with a status variable:

\[
  y \in \{\text{present}, \text{absent}\}.
\]

For target-present examples, the model should preserve useful search behavior and avoid rejecting visible targets. For target-absent examples, the model should avoid hallucinating a target and instead report absence. We evaluate this as a binary present/absent decision, emphasizing absent precision, absent recall, absent F1, and two error types:

\begin{itemize}
  \item \textbf{Absent false-present}: the target is absent, but the model predicts present. This is the forced-choice failure mode.
  \item \textbf{Present false-absent}: the target is present, but the verifier rejects it. This is the main cost of conservative stopping.
\end{itemize}

This formulation deliberately separates target-presence verification from scanpath generation. The current method is post-hoc: it does not claim to improve the generated scanpath. Instead, it asks whether the scanpath and UI evidence contain enough information to decide when a target-present model should stop and reject the query.
